
\documentclass{article}
\usepackage{colortbl}
\usepackage{makecell}
\usepackage{multirow}
\usepackage{supertabular}

\begin{document}

\newcounter{utterance}

\centering \large Interaction Transcript for game `cladder', experiment `full\_v1.5\_default', episode 3593 with Qwen3.5{-}27B.
\vspace{24pt}

{ \footnotesize  \setcounter{utterance}{1}
\setlength{\tabcolsep}{0pt}
\begin{supertabular}{c@{$\;$}|p{.15\linewidth}@{}p{.15\linewidth}p{.15\linewidth}p{.15\linewidth}p{.15\linewidth}p{.15\linewidth}}
    \# & \multicolumn{2}{c}{Player} && \multicolumn{2}{c}{Game Master} \\
    \hline

    \theutterance \stepcounter{utterance}  
    & & & \multicolumn{4}{p{0.6\linewidth}}{
        \cellcolor[rgb]{0.9,0.9,0.9}{
            \makecell[{{p{\linewidth}}}]{
                \texttt{\tiny{[P1$\langle$GM]}}
                \texttt{You are an expert in causal inference. The following question is not a typical commonsense query, but rather a meticulously designed question created by a professor specializing in causal inference, intended to assess the students' mastery of the course content.} \\
\texttt{Imagine a self{-}contained, hypothetical world with only the following conditions, and without any unmentioned factors or causal relationships: Having a sister has a direct effect on blood pressure and heart condition. Blood pressure has a direct effect on heart condition. Method 1: We look at how having a sister correlates with heart condition case by case according to blood pressure. Method 2: We look directly at how having a sister correlates with heart condition in general. To understand how having a sister affects heart condition, is it more correct to use the Method 1 than Method 2?} \\
\texttt{Start your answer with 'Yes' or 'No', followed by additional reasoning or evidence to support your explanation.} \\
            }
        }
    }
     \\ \\

    \theutterance \stepcounter{utterance}  
    & \multicolumn{4}{p{0.6\linewidth}}{
        \cellcolor[rgb]{0.561,0.737,0.561}{
            \makecell[{{p{\linewidth}}}]{
                \texttt{\tiny{[P1$\rangle$GM]}}
                \texttt{No} \\
\\ 
\texttt{In this hypothetical scenario, "Blood Pressure" acts as a **mediator** (or intermediate variable) in the causal pathway from "Having a sister" to "Heart condition." The problem states that having a sister directly affects blood pressure, and blood pressure directly affects heart condition. This creates a causal chain: Having a sister \$\textbackslash{}rightarrow\$ Blood Pressure \$\textbackslash{}rightarrow\$ Heart condition. Additionally, having a sister has a *direct* effect on heart condition as well.} \\
\\ 
\texttt{**Method 2** (looking at the general correlation) estimates the **Total Effect** of having a sister on heart condition. This total effect includes both the direct path (Sister \$\textbackslash{}rightarrow\$ Heart) and the indirect path mediated through blood pressure (Sister \$\textbackslash{}rightarrow\$ BP \$\textbackslash{}rightarrow\$ Heart). Since the goal is to understand "how having a sister affects heart condition" in its entirety, capturing the full magnitude of the impact is usually the standard objective.} \\
\\ 
\texttt{**Method 1** (stratifying or conditioning on blood pressure) estimates the **Direct Effect** only. By conditioning on the mediator (blood pressure), you effectively "block" the indirect pathway. In causal inference, conditioning on a mediator is generally inappropriate if the goal is to estimate the total causal effect, as it leads to an underestimation of the treatment's overall impact by artificially removing the mechanism through which part of the effect operates. This is often referred to as "over{-}adjustment" or "blocking the path."} \\
\\ 
\texttt{Therefore, Method 1 is not more correct than Method 2 for understanding the overall effect; in fact, Method 1 provides a biased estimate of the total effect by excluding the indirect influence transmitted through blood pressure. Method 2 is the correct approach for determining the total causal effect.} \\
            }
        }
    }
    & & \\ \\

    \theutterance \stepcounter{utterance}  
    & & & \multicolumn{2}{p{0.3\linewidth}}{
        \cellcolor[rgb]{0.9,0.9,0.9}{
            \makecell[{{p{\linewidth}}}]{
                \texttt{\tiny{[GM$|$GM]}}
                \texttt{No} \\
            }
        }
    }
    & & \\ \\

    \theutterance \stepcounter{utterance}  
    & & & \multicolumn{2}{p{0.3\linewidth}}{
        \cellcolor[rgb]{0.9,0.9,0.9}{
            \makecell[{{p{\linewidth}}}]{
                \texttt{\tiny{[GM$|$GM]}}
                \texttt{no} \\
            }
        }
    }
    & & \\ \\

    \theutterance \stepcounter{utterance}  
    & & & \multicolumn{2}{p{0.3\linewidth}}{
        \cellcolor[rgb]{0.9,0.9,0.9}{
            \makecell[{{p{\linewidth}}}]{
                \texttt{\tiny{[GM$|$GM]}}
                \texttt{game\_result = WIN} \\
            }
        }
    }
    & & \\ \\

\end{supertabular}
}

\end{document}
